% Copyright (c) 2026 Geovana Neves. Licensed under CC BY 4.0 (https://creativecommons.org/licenses/by/4.0/). See LICENSE-AND-AUTHORSHIP.md.
%
% 05-pproc-definitions/main.tex
%
% Guide 5 of the pyflightstream guides: Post-processing definitions. Built by ../build-all.ps1
% (or build-all.sh) with this folder and ../shared/ as include directories.
%
% No em dashes and no en dashes. No underscores in frame or section titles
% outside \code: beamer writes titles to .nav and .toc.

\documentclass[aspectratio=169,11pt]{beamer}

\input{preamble-slides.tex}
\input{hooks-slides.tex}
\ftsinput{boxes}
\ftsinput{hooks-contract}

\usepackage[nohyphen]{underscore}

\input{info.tex}

\begin{document}

\guidetitleframe{5}{Post-processing definitions}{Post-processing definitions}
  {What every product, reduction and window means}
  {Source: the package page ``Post-processing definitions'' [1].}

\begin{frame}{The definition of record}
\begin{itemize}
  \item \code{docs/post-processing-definitions.md} defines every
        post-processing product the package writes.
  \item Each definition states the requirement a product is built to.
  \item Where code and page disagree, the page is right and the code is a defect.
\end{itemize}
\begin{takeaway}
Read the definitions before changing a reduction, an averaging window or a
product's columns. The implementation is not its own specification.
\end{takeaway}
\end{frame}

\begin{frame}{Contents}
\small
\begin{multicols}{2}
\tableofcontents
\end{multicols}
\end{frame}

% Copyright (c) 2026 Geovana Neves. Licensed under CC BY 4.0 (https://creativecommons.org/licenses/by/4.0/). See LICENSE-AND-AUTHORSHIP.md.
%
\section{The vocabulary}

\begin{frame}{Turns, passages and windows}
\relax
{\ftstablesize
\begin{tabular}{@{}p{0.24\linewidth}p{0.69\linewidth}@{}}
\toprule
\textbf{Term} & \textbf{Meaning} \\
\midrule
Revolution & One full turn of a rotor:
  $60/\mathrm{rpm}$ seconds, $60/(\mathrm{rpm}\,\mathrm{dt})$ solver steps \\
Blade passage & One revolution divided by the blade count \\
Azimuthal position & Where a blade is in its turn, 0 to 360 degrees \\
Window & An inclusive, 1-based range of solver steps that a reduction averages over \\
\bottomrule
\end{tabular}}
\end{frame}

\begin{frame}{Rotor names and reference points}
\relax
{\ftstablesize
\begin{tabular}{@{}p{0.24\linewidth}p{0.69\linewidth}@{}}
\toprule
\textbf{Term} & \textbf{Meaning} \\
\midrule
\code{SMRP} & A rotor's own static moment reference point, \code{<ALIAS>\_SMRP} \\
\code{MRP} & The global moment reference point of the aircraft \\
Alias & The name of a rotor, and of the integration group built from its families \\
\bottomrule
\end{tabular}}
\end{frame}

% Copyright (c) 2026 Geovana Neves. Licensed under CC BY 4.0 (https://creativecommons.org/licenses/by/4.0/). See LICENSE-AND-AUTHORSHIP.md.
%
\section{What every product states}

\begin{frame}{The condition block: angles and velocities}
\small Every table the post stage composes states its condition in one block,
in the order shown across these two frames.
\par\vspace{3mm}
{\ftstablesize
\begin{tabular}{@{}p{0.22\linewidth}p{0.10\linewidth}p{0.59\linewidth}@{}}
\toprule
\textbf{Column} & \textbf{Unit} & \textbf{Meaning} \\
\midrule
\code{ALPHA}, \code{BETA} & deg & Angles the solver reports it ran at, as the row wrote them \\
\code{MACH} & - & Mach number of that point \\
\code{RE} & millions & Reynolds number the solver reports \\
\code{VINF} & m/s & Free-stream velocity the solver reports \\
\code{VREF} & m/s & Solver reference velocity, used to normalise a coefficient \\
\bottomrule
\end{tabular}}
\end{frame}

\begin{frame}{The condition block: atmosphere and references}
\relax
{\ftstablesize
\begin{tabular}{@{}p{0.22\linewidth}p{0.10\linewidth}p{0.59\linewidth}@{}}
\toprule
\textbf{Column} & \textbf{Unit} & \textbf{Meaning, continuing the block} \\
\midrule
\code{ALT} & ft & Altitude the row states; \code{NA} where it states none \\
\code{RHO} & kg/m$^3$ & Air density the run resolved for that point \\
\code{TEMP} & K & Air temperature the run resolved for that point \\
\code{MU} & Pa s & Dynamic viscosity, written in scientific notation \\
\code{J} & - & Advance ratio the row requested; \code{NA} on a row that turns no rotor \\
\code{J\_CLOCK} & - & Advance ratio the clock rotor ran at, $V/(nD)$; \code{NA} where the record or the reference does not say \\
\code{RPM\_CLOCK} & rev/min & Speed the clock rotor turned at, with its hand; \code{NA} with no clock \\
\code{SREF}, \code{CREF}, \code{BREF} & m$^2$, m, m & Reference area, chord and span \\
\bottomrule
\end{tabular}}
\end{frame}

\begin{frame}{A coefficient needs its normalisation}
\small A coefficient is a force divided by
\[
  \frac{1}{2}\rho V^2 S.
\]
\begin{itemize}
  \item Area alone does not let a reader recover a force or compare campaigns.
        Density and the velocity used in the division must be stated too.
  \item \code{ALT} cannot stand in for density: a row may pin density directly.
\end{itemize}
\end{frame}

\begin{frame}{Which velocity normalises which product}
\begin{itemize}\small
  \item The solver uses \code{VREF}. The steady polar's twenty-four
        coefficients are the solver's own and use \code{VREF}.
  \item The plots table, every reduction of it and the unsteady polar are
        rescaled to the free stream by
        $\left(\mathrm{VREF}/\mathrm{VINF}\right)^2$ and use \code{VINF}.
  \item The rotor table recovers Newtons with \code{VREF}, undoing the
        solver's division. Its coefficients use $\rho n^2 D^4$ and carry
        no velocity.
\end{itemize}
\begin{takeaway}
Both velocities are stated. When they differ, post warns and names the point.
\end{takeaway}
\end{frame}

\begin{frame}{Reference area and chord must match the export}
\begin{itemize}\small
  \item The solver divides by the area and length in the project it opened.
        The loads export prints both.
  \item The package sets neither. It states the reference artifact's
        \code{SREF} and \code{CREF}, and checks them against the export.
  \item A difference beyond the export's printed precision is a
        \textbf{warning} in \code{post.log} naming both numbers, and every
        computable product is still written;
        \code{--check-frozen} refuses the simulation's products instead.
\end{itemize}
\begin{takeaway}
A reference beside a coefficient must be the reference that divided it.
\end{takeaway}
\end{frame}

\begin{frame}{Where the condition block appears}
\begin{itemize}
  \item The steady polar already carries \code{ALPHA}, \code{BETA},
        \code{MACH} and \code{RE} among its twenty-four columns. It adds
        the rest beside them.
  \item \code{probes/<point>\_plots.csv} states \textbf{no condition}.
        It keeps the export's own header.
  \item Reductions read every column of that plots table back as a plotted
        quantity.
\end{itemize}
\end{frame}

% Copyright (c) 2026 Geovana Neves. Licensed under CC BY 4.0 (https://creativecommons.org/licenses/by/4.0/). See LICENSE-AND-AUTHORSHIP.md.
%
\section{The axes of a steady polar}

\begin{frame}{One exported force and moment, three axis systems}
\small The loads export gives one force and one moment per surface in the
geometry frame: \textbf{x aft, y right, z up}. Sum over the group's surfaces,
then turn that pair.
\par\vspace{2mm}
{\ftstablesize
\begin{tabular}{@{}p{0.34\linewidth}p{0.21\linewidth}p{0.36\linewidth}@{}}
\toprule
\textbf{Columns} & \textbf{Axes} & \textbf{How} \\
\midrule
\code{CDB, CYB, CLB, CRB25, CMB25, CNB25} & Body: forward, right, down &
Half a turn about y. $\mathrm{CDB}=\mathrm{Cx}$ and $\mathrm{CLB}=\mathrm{Cz}$ of the export \\
\code{CDS ... CNS} & Stability & Body axes turned by $-\alpha_s$ about y \\
\code{CDW ... CNW} & Wind & Stability axes turned by $\beta_w$ about z \\
\bottomrule
\end{tabular}}
\end{frame}

\begin{frame}{Force signs and moment normalisation}
\begin{itemize}
  \item Drag opposes $+x$ and lift opposes $+z$ in each system.
        Side force keeps its sign.
  \item Turn the moment as \textbf{one vector in one length}, then normalise.
  \item Divide \code{CR} and \code{CN} by the span, and \code{CM} by the chord.
  \item Under sideslip, this moves pitch into roll by the ratio of chord to span.
\end{itemize}
\end{frame}

\begin{frame}{The rotation angles come from the flown velocity}
\small The solver turns sideslip about body z first, then incidence. It flies
\[
 (u,v,w)=V(\cos a\cos b,\ \cos a\sin b,\ \sin a),
 \qquad
 \alpha_s=\operatorname{atan2}(w,u),\quad
 \beta_w=\arcsin(v/V).
\]
\begin{itemize}
  \item These angles equal the written \code{ALPHA} and \code{BETA} whenever
        either is zero. Otherwise they differ at second order.
  \item At \code{ALPHA 4, BETA 2}, they are $4.0024^\circ$ and $1.9951^\circ$.
  \item The \code{ALPHA} and \code{BETA} columns keep what the row wrote.
\end{itemize}
\end{frame}

\begin{frame}{Wind-axis drag and the solver's integrals}
\[
  \mathrm{CDW}=\mathrm{CDi}+\mathrm{CDo}.
\]
\begin{itemize}
  \item Wind-axis drag from the export's own vector equals the drag the
        solver integrates, including under sideslip.
  \item Recorded exports confirm this to their printed precision.
  \item \code{CD0} and \code{CDI} state the solver's two integrals.
\end{itemize}
\end{frame}

\begin{frame}{Wind-axis lift comes from the exported vector}
\begin{itemize}\small
  \item \code{CLW} is the vector's wind-axis lift, not the export's \code{CL}.
  \item On 27 of 28 recorded lifting exports with lift above $0.05$,
        the printed \code{CL} is 0.10 to 0.25 per cent higher. One is
        0.71 per cent higher. The cause is unknown.
  \item Every polar axis column uses one source, keeping \code{CDB} and
        \code{CLB} consistent with the export's \code{Cx} and \code{Cz}.
  \item Older tables used the solver's \code{CL} directly, so their
        \code{CLS} and \code{CLW} are higher by those amounts.
\end{itemize}
\vspace{2mm}
{\scriptsize Recorded evidence:
\code{tests/tier1\_offline/fixtures/recorded\_total\_rows.csv}.}
\end{frame}

\begin{frame}{An induced drag the solver did not compute is \code{NA}}
\begin{itemize}\small
  \item A boundary on the vorticity induced-drag list
        (\code{SET\_VORTICITY\_DRAG\_BOUNDARIES}) without a defined trailing
        edge is not computed, and the export prints its \code{CDi} as zero
        [2, p.202].
  \item When the point's run record puts a surface on the list and its printed
        \code{CDi} is exactly zero, the group's \code{CDI} is \code{NA}, and so
        is every axis column the export's x force reaches at that point's
        angles: \code{CDB}, \code{CDS}, \code{CDW}; \code{CLS} and \code{CLW}
        off zero incidence; \code{CYW} under sideslip.
  \item The moments, \code{CYB}, \code{CLB}, \code{CYS} and \code{CD0} keep
        their numbers; the fixed-width custom polar writes \code{nan}.
        \code{post.log} names the declined surfaces.
  \item A trailing-edged surface whose induced drag rounds to zero at the
        printed precision reads \code{NA} too; \code{significant\_digits}
        narrows that band.
\end{itemize}
\end{frame}

% Copyright (c) 2026 Geovana Neves. Licensed under CC BY 4.0 (https://creativecommons.org/licenses/by/4.0/). See LICENSE-AND-AUTHORSHIP.md.
%
\section{The sections table}

\begin{frame}{One export contains every declared distribution}
\begin{itemize}
  \item \code{sections/<point>\_sections.csv} is one solver export.
  \item It holds every distribution the pproc declares: the wing in
        \code{XZ}, each blade in its own frame.
  \item The distributions follow one another with no marker between them.
  \item Leading columns identify which distribution each row belongs to.
\end{itemize}
\end{frame}

\begin{frame}{The sections table: step and distribution identity}
\relax
{\ftstablesize
\begin{tabular}{@{}p{0.17\linewidth}p{0.76\linewidth}@{}}
\toprule
\textbf{Column} & \textbf{Meaning} \\
\midrule
\code{STEP} & Unsteady: the run's last time step, from the run record.
The export is written when the march ends; its own header counts inner
iterations, not time steps. Steady: the solver iteration the export states.
One name across every table \\
\code{FAMILY} & Geometry families of the row's distribution, joined by \code{+} \\
\code{PLANE} & Cutting plane of that distribution \\
\bottomrule
\end{tabular}}
\end{frame}

\begin{frame}{The sections table: rotor identity}
\relax
{\ftstablesize
\begin{tabular}{@{}p{0.17\linewidth}p{0.76\linewidth}@{}}
\toprule
\textbf{Column} & \textbf{Meaning, after \code{PLANE}} \\
\midrule
\code{ROTOR} & Rotor whose blades those families are;
\code{NA} for a surface no rotor owns \\
\code{AZIMUTH} & Where blade one of that rotor is at \code{STEP}, in degrees,
wrapped to one turn; \code{NA} without a rotor \\
\bottomrule
\end{tabular}}
\end{frame}

\begin{frame}{Each rotor has its own azimuth clock}
\small
\[
\mathrm{AZIMUTH}=\left(
  \mathrm{blade1.azimuth\_deg}
  +\mathrm{sense}\,\mathrm{STEP}\,
   \frac{360}{\mathrm{steps\_per\_revolution}}
\right)\bmod 360.
\]
\begin{itemize}
  \item The datum, the sense of rotation and the steps per revolution all
        come from \textbf{that rotor}.
  \item Sense is the sign of the rotor's speed.
  \item Two rotors at two speeds have two azimuths at the same step.
        A wing has none.
\end{itemize}
\end{frame}

\begin{frame}{Distribution identity comes from the run record}
\begin{itemize}
  \item The run records which distribution is which. The script states
        surfaces by index; post cannot recover their names from that alone.
  \item An older record writes \code{NA} in all four identity
        fields: \code{FAMILY}, \code{PLANE}, \code{ROTOR}, \code{AZIMUTH}.
  \item The same rule applies when the recorded blocks do not add up to
        the number of exported rows.
\end{itemize}
\end{frame}

\begin{frame}{The sections table is one instant}
\begin{itemize}
  \item At an unsteady point, the distribution is at \code{STEP}.
        It is not an average over the window.
  \item Its \code{products.json} entry states \code{"kind": "instant"}.
  \item The history is \code{series/<point>\_sections\_series.csv},
        with the same identity on every row.
\end{itemize}
\end{frame}

\begin{frame}{One loads file and one Cp file per distribution}
\begin{itemize}\small
  \item Each \code{[[sections.distributions]]} entry gets
        \code{sections/<point>\_sloads\_<name>.csv} and
        \code{sections/<point>\_cp\_<name>.csv} when its export is available.
  \item The name preserves the alias or joins the declared families with
        \code{-}. Filename-invalid characters become \code{\_}; collisions
        receive the entry's 1-based position suffix.
  \item All planes and expanded blade blocks of that entry share its file.
        Cp now has a table, rather than only a manifest listing.
  \item With per-step exports enabled, every exported step appears in ascending
        \code{STEP} order. Without them, the file holds the end-of-run export.
  \item The combined sections table and combined sections series remain available.
\end{itemize}
\end{frame}

\begin{frame}{Distribution columns and recorded identity}
\begin{itemize}\small
  \item Rows start with \code{STEP, time\_s, FAMILY, PLANE, ROTOR, AZIMUTH},
        then the condition block and export columns. Cp also carries
        \code{SECTION}, the export's 1-based cross-section index.
  \item Missing values are \code{NA}. The manifest records the distribution,
        original families and \code{steps\_tabled}.
  \item A new run records each block's distribution entry. A record can
        be split only when its recorded layout and pproc identify the blocks
        unambiguously by families, plane, frame and count.
  \item Missing or ambiguous identity, mismatched counts, malformed exports
        and unavailable steps are named skips in \code{products.json}.
        Editing today's pproc cannot supply missing recorded identity.
\end{itemize}
\end{frame}

% Copyright (c) 2026 Geovana Neves. Licensed under CC BY 4.0 (https://creativecommons.org/licenses/by/4.0/). See LICENSE-AND-AUTHORSHIP.md.
%
\section{\code{time\_average}}

\begin{frame}{One average of the whole point}
\begin{itemize}
  \item \code{time\_average} is one average of the whole point over one window.
  \item An unsteady POLAR row is built from this average.
\end{itemize}
\begin{takeaway}
One window per point, not one per rotor: a point has one history.
\end{takeaway}
\end{frame}

% Copyright (c) 2026 Geovana Neves. Licensed under CC BY 4.0 (https://creativecommons.org/licenses/by/4.0/). See LICENSE-AND-AUTHORSHIP.md.
%
\section{\code{per\_blade}}

\begin{frame}{One shared window, one row per blade}
\begin{itemize}\small
  \item \code{probes/<point>\_per\_blade\_<ALIAS>.csv} averages each blade
        over the last part of the run: the matrix row's averaging window.
  \item Every blade shares that window. Each row states that blade's start
        and end azimuth.
  \item The window is the one the unsteady plots and POLAR use, from
        \code{LAST\_REVS\_AVG} or \code{LAST\_ITERS\_AVG} on the matrix row.
  \item There is one row per blade.
\end{itemize}
\end{frame}

\begin{frame}{The same time interval makes blades comparable}
\begin{itemize}
  \item Each blade starts at a different azimuth. Its azimuth columns
        state that difference explicitly.
  \item A separate passage window for each blade would also sample each
        blade at a different time, mixing azimuth and history differences.
  \item Averaging from step one mixes the transient with the answer.
\end{itemize}
\begin{takeaway}
The last shared window keeps the sampling time the same for every blade.
\end{takeaway}
\end{frame}

\begin{frame}{The leading columns of each blade row}
\small Each row leads with these fields, in this order:
\begin{enumerate}
  \item \code{REDUCTION}, \code{ROTOR}, \code{BLADE}, \code{FAMILY}.
  \item The window: \code{FIRST\_STEP}, \code{LAST\_STEP}, \code{STEPS}.
  \item \code{AZIMUTH\_START}, \code{AZIMUTH\_END}.
  \item The condition block, then the moment point.
\end{enumerate}
\end{frame}

\begin{frame}{Blade plot columns share one heading}
\begin{itemize}\small
  \item A blade's columns are plots named for its family.
        A plot is \code{<parameter>\_<group>}; a group cut per blade is
        named for the blade's family.
  \item Blade one's \code{CL\_MRP\_Blade1} and \code{FX\_LOCAL\_Blade1}
        become \code{CL\_MRP} and \code{FX\_LOCAL} in its row.
        Removing the family makes the blades line up under one heading.
  \item The pproc requests these with \code{families = "each"} or
        \code{frame = "LOCAL\_AXIS"} over the rotor.
\end{itemize}
\end{frame}

\begin{frame}{Start and end azimuth belong to that blade}
\small At the first and last step of the shared window, use
\[
\begin{split}
\mathrm{azimuth}=\biggl(&\mathrm{blade1.azimuth\_deg}
 +\mathrm{blade\ position}\,\frac{360}{\mathrm{blades}}\\
 &+\mathrm{sense}\,\mathrm{step}\,
 \frac{360}{\mathrm{steps\_per\_revolution}}\biggr)\bmod 360.
\end{split}
\]
\begin{itemize}
  \item \code{blades} is the rotor's count. A periodic sector carrying
        two blades of four spaces them a quarter turn apart.
  \item Without the rotor's clock, the azimuths are \code{NA}.
        The averages are still written.
\end{itemize}
\end{frame}

\begin{frame}{The reference identifies the blade families}
\begin{itemize}
  \item The rotor block's \code{families\_blades} in the reference artifact
        identifies its blades.
  \item The run also records these families.
  \item A row using flat rotor keys and citing no block gets no per-blade
        table. \code{products.json} states why.
\end{itemize}
\end{frame}

% Copyright (c) 2026 Geovana Neves. Licensed under CC BY 4.0 (https://creativecommons.org/licenses/by/4.0/). See LICENSE-AND-AUTHORSHIP.md.
%
\section{\code{phase\_locked}}

\begin{frame}{A mean at a fixed azimuth across revolutions}
\begin{enumerate}\small
  \item Take the last \code{last\_revolutions\_avg} revolutions.
  \item At each azimuth, average the samples at that same position across
        those revolutions.
  \item Tabulate the result by azimuthal position, 0 to 360 degrees.
\end{enumerate}
\begin{takeaway}
With three revolutions, three datapoints enter each azimuth's mean.
The row is an azimuthal position, not a passage, revolution or blade.
\end{takeaway}
\end{frame}

\begin{frame}{Blade and alias means in one file}
\begin{itemize}
  \item Do the fixed-azimuth average for each blade on its own, about that
        rotor's \code{SMRP}.
  \item Then do it for each alias, about the global \code{MRP}.
  \item Put all of it in one file with columns suffixed
        \code{\_\{alias or blade\_name\}\_\{SMRP or MRP\}}.
\end{itemize}
\end{frame}

\begin{frame}{The azimuthal table's leading columns}
\small \code{probes/<point>\_phase\_locked[\_<ALIAS>].csv} leads with:
\begin{enumerate}
  \item \code{REDUCTION}, \code{ROTOR}, \code{AZIMUTH}, \code{STEP},
        \code{REVOLUTIONS}.
  \item The steps the revolutions span:
        \code{FIRST\_STEP}, \code{LAST\_STEP}, \code{STEPS}.
  \item The condition block and the moment point.
  \item The plotted columns, under the names the export prints.
\end{enumerate}
\end{frame}

\begin{frame}{Rows use the rotor's last revolution}
\begin{itemize}
  \item There is one row per solver step of the rotor's last revolution,
        sorted by azimuth from 0 towards 360.
  \item \code{AZIMUTH} is where blade one is, using the sections-table formula.
  \item \code{STEP} is the step of the last revolution that azimuth falls on.
  \item \code{REVOLUTIONS} is the number of samples entering the mean.
        It equals \code{last\_revolutions\_avg} when that is a whole number.
\end{itemize}
\end{frame}

\begin{frame}{Blade columns align by each blade's own azimuth}
\begin{itemize}\small
  \item A column ending in a blade family of the rotor is sampled where
        \textbf{that blade} is at the row's azimuth.
  \item The blade sits its position times $360/\mathrm{blades}$ after blade one.
        This offset makes two blades line up azimuth for azimuth.
  \item Every other column, including a rotor total or the aircraft, is
        tabulated by blade one's azimuth.
\end{itemize}
\end{frame}

\begin{frame}{Azimuthal suffixes come from the plots}
\begin{itemize}\small
  \item A plot is \code{<parameter>\_<group>}. The pproc names the group.
  \item The suffix \code{\_\{alias or blade\_name\}\_\{SMRP or MRP\}}
        is what a group named \code{PUSHER\_SMRP} or cut per blade prints.
  \item The package does not rename these columns.
\end{itemize}
\end{frame}

\begin{frame}{Samples between history steps are read linearly}
\begin{itemize}
  \item One revolution earlier falls a whole number of steps earlier only
        when \code{steps\_per\_revolution} is whole.
  \item A blade offset falls on a whole step only when the steps per
        revolution divide by the blade count.
  \item Where both hold, every sample is a history row and nothing is
        interpolated. Between steps, the history is read linearly.
\end{itemize}
\end{frame}

\begin{frame}{The azimuthal mean needs a rotor clock and depth}
\begin{itemize}
  \item Without blade one's datum or the rotor's signed speed, an azimuth
        cannot be stated. The table is a named skip.
  \item Declare the rotor in the reference and have the matrix row cite it.
  \item A depth below one revolution is a named skip too.
\end{itemize}
\end{frame}

\begin{frame}[fragile]{The pproc declares the gate and averaging depth}
\begin{lstlisting}[style=ftsbase]
[phase_locked]
min_revolutions      = 4.0
last_revolutions_avg = 2.0
\end{lstlisting}
\begin{itemize}\small
  \item \code{min\_revolutions} is the minimum \textbf{total} revolutions
        required for the table to exist. Equal or greater generates it.
  \item Compare the revolutions the matrix row states it turns,
        not the length of the exported window.
  \item \code{last\_revolutions\_avg} is the depth of each azimuthal mean
        and may not exceed \code{min\_revolutions}.
\end{itemize}
\end{frame}

\begin{frame}{The gate and depth can change at post}
\begin{itemize}
  \item The declared \code{[phase\_locked]} table binds.
  \item \code{pyfs-matrix post} reads the current pproc again.
  \item Editing the gate or depth needs no solver re-run.
\end{itemize}
\end{frame}

\begin{frame}{A short run loses only the azimuthal reduction}
\begin{itemize}
  \item A row turning fewer total revolutions than the minimum skips this
        reduction.
  \item Its POLAR, per-blade table and every other product are written as usual.
\end{itemize}
\begin{takeaway}
Missing the threshold skips the reduction; it does not refuse the run's products.
\end{takeaway}
\end{frame}

\begin{frame}{An undeclared table keeps the passage series}
\begin{itemize}
  \item A pproc without \code{[phase\_locked]} still gets the averaging
        window cut into consecutive blade passages, one row per passage.
  \item This is the passage series written under this name.
  \item Declaring the table selects the fixed-azimuth mean instead.
        Leaving it absent preserves the passage series.
\end{itemize}
\begin{takeaway}
Absent is not zero: there is no minimum gate on the undeclared passage series.
\end{takeaway}
\end{frame}

% Copyright (c) 2026 Geovana Neves. Licensed under CC BY 4.0 (https://creativecommons.org/licenses/by/4.0/). See LICENSE-AND-AUTHORSHIP.md.
%
\section{The averaging window}

\begin{frame}{The matrix row states the averaging window}
\relax
{\ftstablesize
\begin{tabular}{@{}p{0.26\linewidth}p{0.24\linewidth}p{0.12\linewidth}p{0.28\linewidth}@{}}
\toprule
\textbf{Column} & \textbf{Run type} & \textbf{Required} & \textbf{Unit} \\
\midrule
\code{LAST\_REVS\_AVG} & \code{unsteady\_rotor} & Yes & Last revolutions; accepts a float \\
\code{LAST\_ITERS\_AVG} & \code{unsteady} & Yes & Last iterations \\
\bottomrule
\end{tabular}}
\vspace{3mm}
\small The window belongs on the matrix row beside \code{DELTA\_TIME},
\code{TIME\_ITERATIONS} and \code{RPM}. It does not belong in the pproc:
these quantities define its temporal setup.
\end{frame}

\begin{frame}{Window keys use the exact upper-case spelling}
\begin{itemize}
  \item \code{VAR\_NAMES\_VALUES} keys match the exact spelling.
  \item A lower-case \code{last\_revs\_avg} on a workflow row is refused
        as a key of no run type.
  \item Where that check does not run, the key is not read at all and
        its intended window is not applied.
\end{itemize}
\end{frame}

\begin{frame}{Retired window keys are refused; the refusal names the key to write}
\begin{itemize}\small
  \item \code{WINDOW\_STEPS}, \code{WINDOW\_REVOLUTIONS} and
        \code{WINDOW\_DEGREES} are retired.
  \item A row using one is refused at plan time, naming the replacement,
        even beside its replacement: \code{LAST\_ITERS\_AVG} for steps,
        \code{LAST\_REVS\_AVG} for revolutions, degrees divided by 360.
  \item These keys bind the \code{time\_average} window and the passages
        cut from it. They do not bind the unsteady POLAR window.
\end{itemize}
\end{frame}

\begin{frame}{New plans and older records without current keys}
\begin{itemize}
  \item A new plan with \textbf{no window key} is refused.
  \item A retired \code{WINDOW\_*} key no longer permits planning: it
        is refused with the key to write instead.
  \item An older record without the current keys is averaged over the
        window the run was given, with a warning naming the steps.
\end{itemize}
\end{frame}

\begin{frame}{Retired compatibility forms are refused by name}
\small
Each retired form is refused, naming its replacement:
\code{iteration=} (use \code{step=}),
\code{run.assess\_unsteady\_from\_plots} (use \code{LoadsAssessor}), the three
\code{WINDOW\_*} row keys (\code{LAST\_ITERS\_AVG} or \code{LAST\_REVS\_AVG};
degrees divided by 360) and group-member lists (one alias per group, members in
the reference's \code{[aliases]}). The \code{broken\_commands} reader stays as
plain compatibility with no countdown; existing run records are never rewritten.
\end{frame}

% Copyright (c) 2026 Geovana Neves. Licensed under CC BY 4.0 (https://creativecommons.org/licenses/by/4.0/). See LICENSE-AND-AUTHORSHIP.md.
%
\section{The unsteady POLAR}

\begin{frame}{The unsteady POLAR averages the plots history}
\begin{itemize}\small
  \item The source is the plots history, time-averaged over the row's window.
        It does not read the native coefficient export for that average.
  \item Columns are plot variables under the names the export prints,
        rather than the steady polar's fixed twenty-four columns.
  \item Flight-condition and reference-length columns are added from the workspace.
  \item The native export states only the last time step, one instant of a
        rotor cycle. It still ships as a health check.
\end{itemize}
\end{frame}

\begin{frame}{One file per simulation, one row per point}
\small \code{polars/P<sim>\_<name>\_uns\_avg.csv}
\begin{itemize}
  \item Every file under \code{post/} comes from a sweep.
        The name states that this is an average of unsteady history.
  \item It carries the \code{P} used by every per-point product.
\end{itemize}
\end{frame}

\begin{frame}{Each polar row states its window and moment point}
\begin{enumerate}
  \item \code{FIRST\_STEP}, \code{LAST\_STEP}, \code{STEPS} open the row
        and identify that point's averaging window.
  \item The condition block follows.
  \item Then \code{XMOM}, \code{YMOM}, \code{ZMOM} state the moment point.
\end{enumerate}
\begin{takeaway}
The plots carry moments. A moment needs its point.
\end{takeaway}
\end{frame}

\begin{frame}{The super-file content joins the polar}
\small After the plot columns, the table adds:
\begin{itemize}
  \item The matrix row's cells and the run record's scalars.
  \item Each rotor's speed.
  \item The solver flags.
\end{itemize}
\begin{takeaway}
For an unsteady point, this content is part of the polar, not a second super file.
\end{takeaway}
\end{frame}

\begin{frame}{Each global plot group supplies eighteen axis columns}
\begin{itemize}\small
  \item The axis coefficients follow the plot columns, using the steady
        polar's names in this order:
        \code{CDW .. CNW25}, \code{CDS .. CNS25}, \code{CDB .. CNB25}.
  \item Each name takes the plot group's \textbf{whole} name as a suffix,
        for example \code{CLW\_MRP\_TOTAL}.
  \item A second global-frame group adds its own eighteen columns
        and renames none.
\end{itemize}
\end{frame}

\begin{frame}{The unsteady axes use dimensional global loads}
\begin{itemize}\small
  \item Use \code{FX, FY, FZ, MX, MY, MZ} from a pproc plot group with
        \code{frame = "MRP"}, in Newtons and Newton metres.
  \item Average these components over the row's window.
        Divide by that row's
        $\tfrac12\mathrm{RHO}\,\mathrm{VINF}^2\,\mathrm{SREF}$
        and also by \code{CREF} for moments.
  \item Turn them as the steady-polar axes are turned, with the same
        force signs and final moment normalisation.
  \item A rotor's own frame cannot supply this block: its axes are not
        the geometry's axes.
\end{itemize}
\end{frame}

\begin{frame}{Missing global plots and the axes block}
\begin{itemize}\small
  \item If the pproc has no global-frame group plotting all six components,
        the run adds \code{MRP\_TOTAL} over every boundary, where the run
        has an \code{MRP} frame.
  \item A pproc that already plots them stays unchanged, to the byte.
  \item An older run without those plots has no axes block.
        \code{products.json} records this under \code{polars/<file>\#axes}.
\end{itemize}
\begin{takeaway}
Missing source plots do not produce an axes block filled with \code{NA}.
\end{takeaway}
\end{frame}

\begin{frame}[fragile]{The pproc can declare output column names}
\small The package cannot infer which plot label carries which coefficient.
Undeclared names pass through exactly as the export prints them.
\begin{lstlisting}[style=ftsbase]
[names]
CL_MRP_TOTAL = "CL_TOTAL"
FX_HUB_PUSHER = "FX_PUSHER"
\end{lstlisting}
\small This dictionary maps an exported plot column to the name
a downstream reader wants. Inventing a source label would otherwise risk
an entire column of \code{NA}.
\end{frame}

\begin{frame}{The names dictionary changes selected headings last}
\begin{itemize}\small
  \item It renames the unsteady polar and averaged reductions:
        \code{time\_average} and the passage series.
  \item \code{probes/<point>\_plots.csv} keeps the export's names because
        it is the source table.
  \item Per-blade and azimuthal tables are not renamed: their columns are
        no longer the export's own names.
  \item Axes and \code{[equations]} read export names.
        The dictionary applies last, to the heading alone.
\end{itemize}
\end{frame}

\begin{frame}{The whole dictionary applies, or none of it does}
\begin{itemize}\small
  \item A source column that no plot of the point prints invalidates the
        dictionary for that point.
  \item A destination name already carried by the table does the same.
  \item Every column then retains the export's name, with a warning and
        \code{products.json} entries under \code{polars/<file>\#names}
        and \code{probes/<point>\#names}.
\end{itemize}
\begin{takeaway}
A failed rename never becomes a column of \code{NA}.
\end{takeaway}
\end{frame}

\begin{frame}{Names that are occupied or refused}
\begin{itemize}\small
  \item \code{CL} is taken. The native export's last-step \code{CL},
        \code{CDi}, \code{CDo}, \code{Cx} and the rest already appear
        in the polar's setup content as a health check.
  \item A plotted column cannot take one of those names.
        \code{CL\_TOTAL} is available instead.
  \item Two dictionary entries giving one name, or a name that is not
        one word, are refused when the pproc is read.
\end{itemize}
\end{frame}

% Copyright (c) 2026 Geovana Neves. Licensed under CC BY 4.0 (https://creativecommons.org/licenses/by/4.0/). See LICENSE-AND-AUTHORSHIP.md.
%
\section{Rotor coefficients}

\begin{frame}{Rotor coefficients describe one rotor}
\begin{itemize}
  \item \code{J}, \code{CT}, \code{CQ}, \code{CP}, \code{ETA}, \code{ETAW}:
        one table per rotor, every column suffixed with that rotor's alias.
  \item They are meaningful for one rotor. Summing several rotors mixes
        different diameters and speeds in the normalisation.
\end{itemize}
\[
  J = \frac{V}{|n|\,D}, \quad
  C_T = \frac{T}{\rho n^2 D^4}, \quad
  C_Q = \frac{Q}{\rho n^2 D^5}, \quad
  C_P = 2\pi\,C_Q\,\mathrm{sign}(n), \quad
  \eta = \frac{J\,C_T}{C_P}
\]
\begin{itemize}\small
  \item \(n\) is the signed speed in revolutions per \textbf{second}: its
        magnitude sets \(J\), its sign multiplies \(C_P\). These are the
        classical airscrew coefficients [6, 7].
\end{itemize}
\end{frame}

\begin{frame}{Tip and helical Mach numbers}
\[ \Omega = \frac{2\pi\,\mathrm{RPM}}{60}, \qquad
   M_\mathrm{tip} = \frac{\Omega R}{a}, \qquad
   M_\mathrm{hel} = \frac{\sqrt{V^2 + (\Omega R)^2}}{a} \]
\begin{itemize}\small
  \item Stated for each rotor an \code{unsteady\_rotor} or
        \code{qsteady\_rotor} row turns, for the rotor of a \code{steady} row
        that states \code{RPM}, and for each actuator disc a row names. \(R\) is half the rotor's diameter, or the
        disc's tip radius; \(V\) and \(a\) are the point's own resolved free
        stream and speed of sound. At \(V = 0\), \(M_\mathrm{hel} =
        M_\mathrm{tip}\).
  \item The rotor table ends with \code{MTIP\_<alias>} and
        \code{MHEL\_<alias>}; \code{plan.json} and the run record carry both
        under \code{rotor\_mach}. A point whose condition does not resolve
        reads \code{NA} in both and keeps its row.
  \item Both are geometric: \(V\) is the free stream alone, without the
        rotor's induced velocity, which near hover adds to the speed the tip
        sees. At \(M_\mathrm{hel} \ge 1\) the tip is sonic, beyond the
        linearised compressibility of a panel method, and the plan warns
        once, naming every such point, its rotor and its \(M_\mathrm{hel}\)
        [1, 10].
\end{itemize}
\end{frame}

\begin{frame}{An unsteady rotor row uses the shared history window}
\begin{itemize}
  \item A steady point's row comes from the loads export.
  \item An unsteady point's row averages the plots history over the
        matrix row's window, shared with the unsteady POLAR and reductions.
  \item The loads export of an unsteady run is the last time step,
        one instant of a cycle.
\end{itemize}
\end{frame}

\begin{frame}{The rotor history is a global group over its own families}
\begin{itemize}\small
  \item Use the rotor's six components \code{FX, FY, FZ, MX, MY, MZ}
        in global \code{MRP}, over exactly its general and blade families.
  \item Find them through a \code{[[plots.groups]]} entry with
        \code{frame = "MRP"} and those exact families, whatever its name.
  \item If the pproc plots none, the run adds \code{ROTOR\_<ALIAS>}.
  \item A group in the rotor's own frame cannot supply this history:
        its forces turn with the rotor and its moments are already about the hub.
  \item A group that only shares the alias's name but covers other families
        is not a source either.
\end{itemize}
\end{frame}

\begin{frame}{A windowed rotor row never substitutes an instant}
\begin{itemize}\small
  \item A point whose history does not cover the window, or lacks the required
        columns, is left out and named in \code{products.json}.
  \item The unsteady POLAR applies the same exclusion rule.
  \item The rotor table's manifest entry states \code{source} and \code{window}.
\end{itemize}
\end{frame}

\begin{frame}{\code{ETAW}: the full force vector in wind axes}
\begin{enumerate}
  \item Take the full rotor-frame force vector, $[F_x,F_y,F_z]_{\mathrm{rotor}}$.
  \item Carry it to the airframe body frame with the \textbf{transpose}
        of the rotor-to-body rotation.
  \item Carry it to wind axes with the \textbf{AIAA rotation}, using
        both alpha and beta.
  \item Take its X component, \code{Fx\_W}.
\end{enumerate}
\end{frame}

\begin{frame}{\code{ETAW}: a dimensionless efficiency}
\[
  \mathrm{CTW}=\frac{F_{x,W}}{\rho n^2D^4},
  \qquad
  \mathrm{ETAW}=\frac{J\,\mathrm{CTW}}{\mathrm{CP}}.
\]
\begin{itemize}\small
  \item Wind-axis force replaces thrust, normalised exactly as thrust is.
        \code{CTW} enters where \code{CT} enters the efficiency.
  \item \code{ETAW} reduces to \code{ETA} when the shaft lies along the stream.
  \item No wind-axis force supplied means \code{NA}.
  \item Two vector rotations preserve the off-axis force components.
        The cosine of a scalar angle does not.
\end{itemize}
\end{frame}

\begin{frame}{A static point has only a real zero advance ratio}
\begin{itemize}\small
  \item Every rotor coefficient is \code{NA} at a static point except
        \code{J}, which is \code{0.00000}.
  \item The export gives coefficients normalised by the run's dynamic pressure.
        At $V=0$, that pressure is zero: the real thrust has already been
        divided away before the package sees it.
  \item No rotor coefficient can be recovered from that static export.
  \item A hover figure of merit cannot be offered either: it needs a force
        the run does not state.
\end{itemize}
\end{frame}

% Copyright (c) 2026 Geovana Neves. Licensed under CC BY 4.0 (https://creativecommons.org/licenses/by/4.0/). See LICENSE-AND-AUTHORSHIP.md.
%
% 05-pproc-definitions/text/15-quasi-steady-rotor.tex
%
% The qsteady_rotor workflow's products, read against the merged workflow
% (cases/qsteady.py, post/qsteady.py) and docs/workspace-and-workflows.md.
% The count nP is the blade's own, per revolution; Guide 1 has the rest.

\section{The quasi-steady rotor}

\begin{frame}{A rotor solved steady, and what its products are}
\begin{itemize}\small
  \item A \code{qsteady\_rotor} row holds the blades still and turns the free
        stream about the shaft at the rotor's speed; every solve is steady. It
        is valid only for an \textbf{isolated, axisymmetric} rotor.
  \item \textbf{Sector} (the row states periodic symmetry): one blade and its
        copies, one steady solve. \textbf{Wheel}: every blade meshed; in an
        axial, uniform inflow it is solved once.
  \item A wheel at incidence, or in a custom inflow, sees a different flow at
        each blade position. It is solved once per position and the
        positions are averaged: \code{PASSAGE\_POSITIONS: k}, mandatory,
        \(k \ge 1\), at \(\theta_i = i\,(360^\circ/N)/k\), \(i = 0 \ldots
        k-1\), uniform over one blade passage.
  \item Such a point's products are those of an unsteady point: the loads at
        each position (\code{<point>\_qs<i>.txt}), and their average where an
        unsteady point has its window average.
\end{itemize}
\begin{takeaway}Uniform positions over one passage give an unbiased mean of a
signal that repeats every passage [8]. Two positions settle thrust and
torque; the in-plane forces and moments want six or more.\end{takeaway}
\end{frame}

\begin{frame}{The validity parameter, carried by every product}
\begin{columns}[T]
\begin{column}{0.50\textwidth}
\begin{itemize}\small
  \item Each blade station's 1P reduced frequency,
        \[ k = \frac{\Omega\,c}{2\,V_\mathrm{rel}}, \]
        compares the time the flow takes to cross the chord with the period of
        the once-per-revolution change. Below about 0.05 a quasi-steady
        answer holds; above 0.1 the lag of the circulation and the shed wake
        matter [10, 11, 12].
  \item \pyfs{plan} shows, per wheel point, the per cent of the span with
        \(k > 0.1\) and \(k\) minimum, maximum and mean, and warns when that
        per cent is above zero.
\end{itemize}
\end{column}
\begin{column}{0.46\textwidth}
\begin{block}{Where it is written}
\small
\begin{itemize}\small
  \item The datapoint folder: \(k\) minimum, maximum and mean; the span
        fractions above 0.05 and 0.1; the shares of thrust and torque from
        stations above 0.1.
  \item Every post product of the point carries them.
  \item The sections product carries \(k\) at each station.
\end{itemize}
\end{block}
{\scriptsize With \code{--inflow-fft} a custom inflow's harmonics count too:
\(k_\mathrm{eff} = n_{95}\,k\), \(n\)P counted by one blade per revolution
(Guide 1).\par}
\end{column}
\end{columns}
\end{frame}

% Copyright (c) 2026 Geovana Neves. Licensed under CC BY 4.0 (https://creativecommons.org/licenses/by/4.0/). See LICENSE-AND-AUTHORSHIP.md.
%
\section{What the pproc asks the solver for}

\begin{frame}{Force plots: \code{[plots] parameters}}
\relax
{\ftstablesize
\begin{tabular}{@{}p{0.16\linewidth}p{0.25\linewidth}p{0.29\linewidth}p{0.208\linewidth}@{}}
\toprule
\textbf{In the pproc} & \textbf{The solver's PARAMETER} & \textbf{Units the package requests} & \textbf{Plots column} \\
\midrule
\code{CL}  & \code{CL}       & \code{COEFFICIENTS} & \code{CL\_<group>} \\
\code{CDI} & \code{CDI}      & \code{COEFFICIENTS} & \code{CDI\_<group>} \\
\code{CDO} & \code{CDO}      & \code{COEFFICIENTS} & \code{CDO\_<group>} \\
\code{CD}  & \code{CD}       & \code{COEFFICIENTS} & \code{CD\_<group>} \\
\code{FX}  & \code{FORCE\_X}  & \code{NEWTONS} & \code{FX\_<group>} \\
\code{FY}  & \code{FORCE\_Y}  & \code{NEWTONS} & \code{FY\_<group>} \\
\code{FZ}  & \code{FORCE\_Z}  & \code{NEWTONS} & \code{FZ\_<group>} \\
\code{MX}  & \code{MOMENT\_X} & \code{NEWTONS} & \code{MX\_<group>} \\
\code{MY}  & \code{MOMENT\_Y} & \code{NEWTONS} & \code{MY\_<group>} \\
\code{MZ}  & \code{MOMENT\_Z} & \code{NEWTONS} & \code{MZ\_<group>} \\
\bottomrule
\end{tabular}}
\begin{itemize}\scriptsize
  \item Write the short name. The pproc refuses \code{FORCE\_X}.
        Omitted \code{parameters} means all ten. There is no \code{CY}, \code{CM} or \code{CN} plot.
  \item The solver also offers \code{KILO-NEWTONS}, \code{POUND-FORCE} and
        \code{KILOGRAM-FORCE}. The package fixes one unit token per parameter.
        The plots table rescales coefficient columns by $(\mathrm{VREF}/\mathrm{VINF})^2$.
  \item Global \code{MRP} loads \code{FX .. MZ} feed the unsteady polar and rotor table.
        With \code{MRP} available, the run adds \code{MRP\_TOTAL} if no MRP group plots all six,
        and \code{ROTOR\_<ALIAS>} if none plots all six over that rotor's families.
\end{itemize}
{\tiny\color{SoftGray} Sources under \code{src/pyflightstream/}:
\code{cases/\_\_init\_\_.py}, \code{FORCE\_PLOT\_PARAMETERS}, \code{PlotsSpec};
\code{commands/unsteady\_solver.yaml}, \code{UNSTEADY\_SOLVER\_NEW\_FORCE\_PLOT};
\code{cases/workflows.py}, \code{\_pproc\_plots}, \code{\_plot\_each\_rotors\_own\_history};
\code{post/products.py}, \code{write\_plots\_table}.\par}
\end{frame}

\begin{frame}{Probes: \code{[[probes]] parameters}}
\small The pproc and the solver use the same spelling:
\par\vspace{2mm}
{\ftstablesize
\begin{tabular}{@{}llll@{}}
\code{CP\_FREE} & \code{CP\_REF} & \code{MACH} & \code{VELOCITY} \\
\code{VX} & \code{VY} & \code{VZ} & \code{STATIC\_PRESSURE\_RATIO}
\end{tabular}}
\begin{itemize}\small
  \item On the unsteady fluid-plot path, each parameter becomes one plot per vertex:
        \code{<PARAMETER><n>}, with \code{n} counting across entries.
        The probes table un-pivots these columns to one row per probe and step.
  \item On a steady run, \code{EXPORT\_PROBE\_POINTS} takes no variable list.
        Its columns are the solver's own.
  \item An entry with empty \code{parameters} emits nothing on either run type.
\end{itemize}
\small These six boundary-layer parameters are accepted on builds
26.122, 26.123 and 26.124. An unsupported parameter is refused by name and build:
\par\vspace{2mm}
{\ftstablesize
\begin{tabular}{@{}ll@{}}
\code{BL\_MOMENTUM\_THICKNESS} & \code{BL\_DISPLACEMENT\_THICKNESS} \\
\code{BL\_TOTAL\_THICKNESS} & \code{BL\_SHAPE\_FACTOR} \\
\code{BL\_SKIN\_FRICTION} & \code{BL\_TRANSITION\_MARKER}
\end{tabular}}
\par\vfill
{\tiny\color{SoftGray} Sources under \code{src/pyflightstream/}:
\code{cases/\_\_init\_\_.py}, \code{FLUID\_PLOT\_PARAMETERS}, \code{ProbesSpec};
\code{commands/unsteady\_solver.yaml}, \code{UNSTEADY\_SOLVER\_NEW\_FLUID\_PLOT};
\code{commands/probe\_points.yaml}, \code{EXPORT\_PROBE\_POINTS};
\code{cases/workflows.py}, \code{\_emit\_one\_probe\_table};
\code{post/products.py}, \code{write\_unsteady\_probes\_table}.\par}
\end{frame}

\begin{frame}{The probe source follows the run type}
\begin{itemize}\small
  \item On an unsteady row, every \code{[[probes]]} entry uses fluid plots:
        drawn lines and cited \code{points\_file} profiles alike.
  \item Each vertex and selected parameter shares one plot counter.
        \code{probes/<point>\_probes.csv} holds both forms together,
        one row per probe and time step, from the plots history.
  \item A cited profile is read at plan: a count, then \code{X,Y,Z,TYPE}
        CSV rows. Types 0 and 1 become fixed fluid-plot vertices with the
        entry's frame and scale. Invalid profiles are refused before solving.
  \item The unsteady path no longer imports or exports standard probe points.
        Its default output list has seven files instead of eight.
\end{itemize}
\end{frame}

\begin{frame}{Steady probes and older unsteady records}
\begin{itemize}\small
  \item Steady rows keep the probe-points source for both drawn and cited probes.
        A nonempty \code{parameters} list enables the entry; an empty list disables it.
  \item On a steady row the list does not filter exported variables:
        the solver exports a fixed set. Plan warns with the entry number and frame.
  \item On an unsteady row the list selects the sampled fluid-plot variables.
  \item An older cited profile exported only as an instant has no history.
        Re-post skips it with its reason in \code{products.json}, preserving
        the drawn-line histories that were recorded. A new run supplies
        the missing history.
\end{itemize}
\end{frame}

\begin{frame}[fragile]{What a point exports: \code{[exports]}}
\relax
{\ftstablesize
\begin{tabular}{@{}p{0.17\linewidth}p{0.13\linewidth}p{0.51\linewidth}p{0.098\linewidth}@{}}
\toprule
\textbf{Key} & \textbf{File suffix} & \textbf{Solver command} & \textbf{Run type} \\
\midrule
\code{simulation} & \code{.fsm} & \code{SAVEAS} & Both \\
\code{loads} & \code{.txt} & \code{EXPORT\_SOLVER\_ANALYSIS\_SPREADSHEET} & Both \\
\code{tecplot} & \code{.dat}, \code{.vtk} & \code{EXPORT\_SOLVER\_ANALYSIS\_VTK}; the package writes the \code{.dat} & Both \\
\code{sections} & \code{\_cp.txt} & \code{EXPORT\_ALL\_SURFACE\_SECTIONS} & Both \\
\code{sectional\_loads} & \code{\_sloads.txt} & \code{EXPORT\_SURFACE\_SECTIONAL\_LOADS} & Both \\
\code{probes} & \code{\_probes.txt} & \code{EXPORT\_PROBE\_POINTS} & Steady only \\
\code{plots} & \code{\_plots.txt} & \code{UNSTEADY\_SOLVER\_EXPORT\_PLOTS} & Unsteady only \\
\code{log} & \code{\_log.txt} & \code{EXPORT\_LOG} & Both \\
\bottomrule
\end{tabular}}
\begin{itemize}\small
  \item Keys are \code{true} or \code{false}. The original kinds default on for their run type.
        \code{vtk}, \code{csv} and \code{force\_distributions} default off; the
        steady plot kinds default on (next frame).
\end{itemize}
\end{frame}

\begin{frame}{What a point exports: \code{[exports]} (continued)}
\begin{itemize}\small
  \item \code{loads} is required to judge the run, and \code{simulation}
        is required too: every point leaves its final
        \code{.fsm}. Either one stated \code{false} is refused.
  \item \code{tecplot} exports every surface as VTK and the package writes the
        \code{.dat} from it, per panel, in the reference frame.
  \item \code{singularity\_strength = true}, a top-level pproc key, off by
        default, adds the nodal \code{Singularity\_strength} of a native
        auxiliary \code{*\_native\_tecplot.dat}, matched to the same step's
        VTK by geometry [4]. Off, the point makes no native export, its
        \code{.dat} states the strength \code{NOT\_CARRIED} (so does its
        \code{products.json} entry), and the point is complete without it.
        \pyfs{plan} says per row whether the strength is carried.
\end{itemize}
{\tiny\color{SoftGray} Sources under \code{src/pyflightstream/}:
\code{cases/\_\_init\_\_.py}, \code{EXPORT\_KINDS}, \code{default\_outputs},
\code{\_known\_export\_kinds}; \code{commands/solver\_export.yaml},
\code{EXPORT\_SOLVER\_ANALYSIS\_VTK}; \code{results/surface.py}.\par}
\end{frame}

\begin{frame}{Additional exports: plots, a volume section, force distributions}
\begin{itemize}\small
  \item \textbf{The solver's own plots.} \code{plot\_residuals},
        \code{plot\_loads} (on by default) and \code{plot\_sections\_cp} (on
        where sections are declared) under \code{[exports]}: the plotted series
        as TEXT, one row per iteration or one x/Cp pair per section.
        \textbf{Never a coefficient source}: on 26.124 the plotted pitching
        moment is $-0.897$ where the exported \code{CMy} is $-0.0247$ [3].
        An unsteady row saves the first two as well, once,
        after the march, and the section Cp too, where
        sections are declared.
        \code{false} turns each off.
  \item \textbf{A volume section.} \code{[volume\_section]}: one plane, a
        rectangle (\code{corners\_m}) or an annulus (\code{radii\_m}), in a
        \code{plane} of a \code{frame}, at \code{offset\_m}. \alert{It
        is sampled, not exported by the solver}: probes on a steady
        row, fluid plots on an unsteady or rotor row, written as
        a vertex cloud to \code{post/<matrix>/fields/<point>\_vsec.vtk} or
        \code{.dat} (\code{\_step\_<STEP>} per unsteady STEP). Older records
        keep their native export in the datapoint.
  \item \textbf{The force distribution.} \code{force\_distributions =
        true}: per-panel pressure and viscous force coefficients, once, at the
        end of the run. Off by default.
\end{itemize}
\end{frame}

\begin{frame}[fragile]{VTK and CSV surface exports}
\small \code{vtk\_variables} is a top-level pproc key, before any tables:
\begin{lstlisting}[style=ftsbase]
vtk_variables = ["X", "Y", "Z", "CP_FREESTREAM"]
[exports]
vtk = true
csv = true
\end{lstlisting}
\begin{itemize}\small
  \item Both kinds default off. VTK uses \code{EXPORT\_SOLVER\_ANALYSIS\_VTK}
        and \code{SET\_VTK\_EXPORT\_VARIABLES} when the list is present.
        An absent list requests all variables.
  \item CSV uses \code{EXPORT\_SOLVER\_ANALYSIS\_CSV}: \code{CP-FREESTREAM},
        \code{PASCALS}, all surfaces, solver reference frame. Both exclude the wake.
  \item Commands and variables are checked against the selected build's database.
        Unsupported requests are refused at plan, naming the build.
\end{itemize}
\end{frame}

\begin{frame}[fragile]{The surface averaging window: the package averages}
\begin{lstlisting}[style=ftsbase]
[time_averaging]
last_revs = 1.5 # OR last_iters = 54; exactly one, positive
\end{lstlisting}
\begin{itemize}\small
  \item \alert{The package averages; the solver does not.} The run
        exports the surface at every step of the window, as
        \code{EXPORT\_UNSTEADY\_AFTER\_ITER: <first step>} would, and the post
        averages those exports into \code{surfaces/<point>\_time\_average.dat}:
        the same panel across the steps, every step the same weight, every
        variable. \code{SOLVER\_TIME\_AVERAGING} is never emitted: it hung
        26.124, killed at 240.5 seconds with nothing written [3], and 26.123
        does not carry it.
\end{itemize}
\end{frame}

\begin{frame}{The surface window: time steps, verified on 26.122}
\begin{itemize}\small
  \item Bounds are inclusive, 1-based time steps ending at the run's last step.
        \code{last\_iters} counts steps; \code{last\_revs} uses the same rotor
        clock and rounding as \code{LAST\_REVS\_AVG}. A long window clips at step 1.
  \item Verified on 26.122 [3]: the solver's own average over steps
        7 to 12 is the uniform mean of the run's per-step surfaces over those
        steps, to 1e-13 in Cp, Vx and speed, which is the package's mean. The
        solver keeps CF at its last instant; the package averages it.
  \item Steps that do not share one topology refuse the average; a step of the
        window that was not exported (a clock-stopped run, a continuation) skips
        it. Both by name, never a partial average.
  \item Needs 26.122 or later and \code{tecplot} on. The export header's
        inner-iteration count is not this clock.
\end{itemize}
\end{frame}

\begin{frame}{Surface export provenance and per-step exports}
\begin{itemize}\small
  \item \code{products.json} states the averaged surface \code{kind: average},
        with the recorded window, the steps averaged and each per-step VTK read
        with its sha256; the per-step exports stay \code{kind: instant}.
  \item The window is read from the run record (\code{surface\_average\_window}),
        never from a pproc edited since; changing it requires a new run.
  \item Every \code{.dat} states its source in \code{DATASETAUXDATA}: the VTK,
        its sha256 and a \code{TRANSLATION} record (cell-centred, reference frame).
  \item \code{EXPORT\_UNSTEADY\_AFTER\_REV} or \code{EXPORT\_UNSTEADY\_AFTER\_ITER}
        exports the VTK per step, CSV where asked; each step's \code{.dat} is
        written from its VTK.
  \item An older record keeps the solver's window and its
        meaning. The matrix's plots-reduction window remains a separate choice.
  \item A storage recipe's \code{[[prune\_step\_exports]]} keeps only each
        point's last step; a post that needs a deleted step refuses, naming it
        (Guide 1).
\end{itemize}
\end{frame}

\begin{frame}{Surface formats without a pproc export key}
\begin{itemize}\small
  \item \code{EXPORT\_SOLVER\_ANALYSIS\_PLOAD\_BDF} still has no
        \code{[exports]} key. \code{EXPORT\_SOLVER\_ANALYSIS\_FORCE\_DISTRIBUTIONS}
        has one, \code{force\_distributions}.
  \item Its parameter or payload block cannot be supplied by a setup's
        one-line \code{[[raw]]} entry.
  \item \code{EXPORT\_BL\_VELOCITY\_PROFILE} holds an unattended script
        [3]: a raw line is refused at plan on 26.124.
\end{itemize}
\end{frame}

% Copyright (c) 2026 Geovana Neves. Licensed under CC BY 4.0 (https://creativecommons.org/licenses/by/4.0/). See LICENSE-AND-AUTHORSHIP.md.
%
\section{The additional post}

\begin{frame}[fragile]{A second pproc, extracted from the saved simulation}
\begin{lstlisting}[style=ftsbase]
ADDITIONAL_PPROC: p002        # VAR_NAMES_VALUES of the row
pyfs-matrix post matriz.fs --workspace . --additional-pproc
\end{lstlisting}
\begin{itemize}\small
  \item The row runs byte for byte as without the key; the run record never
        carries it.
  \item Afterwards, each recorded point whose final \code{.fsm} is on disk and
        hashes as its record says is copied into
        \code{datapoints/DP-<point>/additional/<pid>/} and reopened there,
        \textbf{with no solve}: sections cut in the frames the run created,
        sectional loads computed, exports written, file closed. Nothing is
        saved and no probe is written.
  \item Each extraction is recorded in \code{additional.json} beside
        \code{runs.json}, which is never written. The original \code{.fsm} is
        hashed again after the launch.
\end{itemize}
{\tiny\color{SoftGray} Source:
\code{docs/post-processing-definitions.md} [1], The additional post.\par}
\end{frame}

\begin{frame}{What a reopened saved simulation gives back}
\relax
{\ftstablesize
\begin{tabular}{@{}p{0.44\linewidth}p{0.49\linewidth}@{}}
\toprule
\textbf{Export} & \textbf{Reopened, with no solve (26.124) [3]} \\
\midrule
total loads & identical \\
surface solution & identical bytes \\
surface sections, a distribution created after reopening included & identical \\
sectional loads & zero until computed, identical once computed: the file
stores the sections, not their loads, so the extraction computes them every time \\
plots history of an unsteady point & identical \\
probe points off the body & \textbf{not identical}: up to 4 per cent in speed,
0.094 in Cp \\
\bottomrule
\end{tabular}}
\begin{itemize}\small
  \item Only 26.124 was measured: a row on another build is refused at plan.
  \item Probes, a volume section, \code{[plots]}, \code{[time\_averaging]} and
        \code{base\_regions} in an additional pproc are refused: nothing measured
        says a reopened file gives them back.
\end{itemize}
\end{frame}

\begin{frame}{The products of an extraction}
\begin{itemize}\small
  \item Under \code{post/<matrix>/additional/<pid>/}, beside the run's own and
        never over them: the group polars of a steady point, one sections table
        per point, and on an unsteady point the plots tables and reductions.
  \item Every \code{products.json} entry for them carries:
\end{itemize}
{\ftstablesize
\begin{tabular}{@{}p{0.20\linewidth}p{0.72\linewidth}@{}}
\toprule
\textbf{Key} & \textbf{Meaning} \\
\midrule
\code{pproc} & the ADDITIONAL pproc id, not the one the row ran with \\
\code{additional} & always \code{true}; a reader of run products filters on it \\
\code{extraction} & the extractions the file holds, \code{<point run id>/additional/<pid>} \\
\code{derives\_from} & the points those extractions were taken from \\
\bottomrule
\end{tabular}}
\par\vspace{1mm}
\small and no \code{runs}, which names run ids everywhere else in the index.
\end{frame}

\begin{frame}{Which rows, which instant, and when it stops counting}
\begin{itemize}\small
  \item \textbf{Rows.} A reopened sections export carries the run's
        distributions first and the additional pproc's after them. The table
        keeps every row, the new blocks numbered on after the run's and marked
        with the pproc; \code{leading\_sections} counts the run's rows.
  \item \textbf{One instant on an unsteady point.} The saved simulation is the
        LAST instant: the sections table's \code{STEP} is the run's last time
        step, its entry says \code{kind: instant}, and \code{post.log} says so.
        The plots history is the run's own.
  \item \textbf{Only a current extraction has products}: its point is admitted,
        the \code{.fsm} still hashes as the one it opened, and every file it
        wrote is on disk and hashes as recorded. A point run again archives the
        extraction; the stale one is skipped under
        \code{additional/<pid>/runs/<extraction id>}, never under the run's key.
\end{itemize}
\end{frame}

% Copyright (c) 2026 Geovana Neves. Licensed under CC BY 4.0 (https://creativecommons.org/licenses/by/4.0/). See LICENSE-AND-AUTHORSHIP.md.
%
\section{What the package does not judge}

\begin{frame}{Last-step convergence does not judge a settled history}
\begin{itemize}
  \item For an unsteady point, the package asserts whether the numerical
        iterations \textbf{within the last time step} converged.
  \item It does not judge whether the time history has settled.
  \item There is no settle tolerance or convergence criterion over the history,
        and no point fails for one.
\end{itemize}
\begin{takeaway}
The user judges whether the history has settled afterwards, from the history.
\end{takeaway}
\end{frame}

\begin{frame}{Frozen solves and averages}
\begin{itemize}\small
  \item Two consecutive frozen steps make the point \code{FAILED\_DIVERGED}.
        In a frozen step, inner iterations after the first print exactly zero
        velocity residual, and the last prints both residuals exactly zero.
  \item Assessment names the first frozen step and count. Post reads the native
        log again, including for a record previously assessed successful.
  \item An average whose window reaches the first frozen step is skipped with
        its reason in \code{products.json}. An earlier stale product is archived.
  \item Windows wholly before the freeze, raw histories and explicitly instant
        products remain available. A frozen stretch requires a new solve.
\end{itemize}
\end{frame}

\begin{frame}{The post never blocks and always writes its log}
\begin{itemize}\small
  \item Every \code{pyfs-matrix post} writes \code{post.log} beside
        \code{products.json}, even when nothing is wrong; the manifest names it.
        Read it after every post: each named skip and warning states the point,
        the product, the step and what settles it.
  \item By default a frozen point, an unread residual block, a failed status or a
        reference mismatch is a WARNING line, and every computable product is
        written. \code{--check-frozen} refuses the affected averages instead.
  \item The freeze and azimuthal guards judge the exact set of plotted samples the
        reducer reads, blade aliases resolved and every nonzero interpolation
        weight counted; the guard does no arithmetic of its own.
  \item \code{integrate = true} on a \code{[[sections.distributions]]} pproc
        entry appends \code{Strip\_length}, \code{Fx\_int}, \code{Fz\_int} and
        \code{My\_int} to the same sloads file: strips midpoint to midpoint,
        half strips at the ends, the moment about each station's quarter chord.
        Off by default; the file is byte-identical when off.
\end{itemize}
\end{frame}

\begin{frame}{One post, one log, and the log as fields}
\begin{itemize}\small
  \item \code{post.log.json} sits beside \code{post.log}: the same header
        (\code{version}, \code{workspace}, \code{matrix}, \code{time},
        \code{check\_frozen}) and one record per WARNING line, in the same
        order, with \code{point}, \code{product}, \code{message} and
        \code{remedy}. \code{products.json} names it under \code{log\_json}.
  \item Both are written from one list, on a clean, a failed and an
        interrupted post, so they cannot disagree. Parse the JSON, not the text.
  \item A WARNING line names its own point and product:
        \code{WARNING point=camp/sim\_7001/AL-020 product=available-exports: ...}.
  \item Each post collects the package's warnings in its own thread: two posts
        in one process no longer share a log, and a warning from code outside
        the package is not written to it [3].
\end{itemize}
\end{frame}

% Copyright (c) 2026 Geovana Neves. Licensed under CC BY 4.0 (https://creativecommons.org/licenses/by/4.0/). See LICENSE-AND-AUTHORSHIP.md.
%
\section{The token for a value that does not exist}

\begin{frame}{Every product writes \code{NA} for no value}
\begin{itemize}\small
  \item Zero is a real value. Advance ratio at rest is exactly zero,
        so zero cannot stand for a value that does not exist.
  \item A blank cannot distinguish missing data, an inapplicable column
        or a writer that stopped halfway.
  \item The token is defined once in
        \code{pyflightstream.\_tokens.NOT\_APPLICABLE}, below every layer.
  \item Readers still accept \code{-}, the token written in older files.
\end{itemize}
\end{frame}


\begin{frame}{The physics behind the definitions, and where to read it}
\begin{itemize}\small
  \item \textbf{Body, stability and wind axes}, and the angles that turn one
        into the other, are those of flight mechanics [5].
  \item \textbf{The rotor coefficients} \(J\), \(C_T\), \(C_Q\), \(C_P\) and
        \(\eta\) are the classical airscrew coefficients [6, 7].
  \item \textbf{A window average} is the time mean of a record over a finite
        span; its error falls as the span covers more periods of the signal
        [8].
  \item \textbf{The azimuthal (phase-locked) mean} averages the samples that
        share one phase of the revolution, separating the periodic part of a
        signal from the rest [9].
  \item \textbf{Tip and helical Mach numbers}, and the reduced frequency
        that says when a rotor's flow is quasi-steady, are rotor aerodynamics
        [10, 11, 12].
\end{itemize}
\end{frame}

\begin{frame}{References}
\begin{reflist}
  \bibitem{g05r01} \docref{post-processing-definitions}{Post-processing definitions}
  \bibitem{g05r02} Altair Engineering, \emph{FlightStream User Manual}, version 26.12, 2026.
  \bibitem{g05r03} \recordsref
  \bibitem{g05r04} \docref{surface-translation}{Surface exports and native singularity strength}
  \bibitem{g05r05} B. Etkin, L. D. Reid, \emph{Dynamics of Flight: Stability and
        Control}, 3rd ed., Wiley, New York, 1996.
  \bibitem{g05r06} H. Glauert, \emph{The Elements of Aerofoil and Airscrew Theory}, 2nd
        ed., Cambridge University Press, Cambridge, 1947.
  \bibitem{g05r07} B. W. McCormick, \emph{Aerodynamics, Aeronautics, and Flight
        Mechanics}, 2nd ed., Wiley, New York, 1995.
  \bibitem{g05r08} J. S. Bendat, A. G. Piersol, \emph{Random Data: Analysis and
        Measurement Procedures}, 4th ed., Wiley, Hoboken, 2010.
  \bibitem{g05r09} W. C. Reynolds, A. K. M. F. Hussain, ``The mechanics of an organized
        wave in turbulent shear flow. Part 3. Theoretical models and
        comparisons with experiments'', \emph{Journal of Fluid Mechanics},
        54(2), pp. 263-288, 1972.
  \bibitem{g05r10} J. G. Leishman, \emph{Principles of Helicopter Aerodynamics}, 2nd
        ed., Cambridge University Press, Cambridge, 2006.
\end{reflist}
\end{frame}

\begin{frame}{References (continued)}
\begin{reflist}
  \setcounter{enumi}{10}
  \bibitem{g05r11} T. Theodorsen, ``General theory of aerodynamic instability and the
        mechanism of flutter'', NACA Report 496, 1935.
  \bibitem{g05r12} W. R. Sears, ``Some aspects of non-stationary airfoil theory and its
        practical application'', \emph{Journal of the Aeronautical Sciences},
        8(3), pp. 104-108, 1941.
\end{reflist}
\end{frame}

\end{document}
